\section{Bài tập chuyên đề}

\subsection{Củng cố}

\begin{tuluyen}
Viết công thức cấu tạo thu gọn của ethylic alcohol, acetic acid và ethyl acetate. Chỉ ra nhóm \ce{-OH}, nhóm \ce{-COOH} và nhóm ester trong các công thức tương ứng.
\dapso{\ce{CH3CH2OH}; \ce{CH3COOH}; \ce{CH3COOC2H5}.}
\end{tuluyen}

\begin{tuluyen}
Viết phương trình hóa học của các phản ứng sau:
\begin{tasks}(2)
\task ethylic alcohol cháy hoàn toàn;
\task ethylic alcohol tác dụng với Na;
\task acetic acid tác dụng với Mg;
\task acetic acid tác dụng với NaOH;
\task acetic acid tác dụng với CuO;
\task acetic acid tác dụng với \ce{CaCO3};
\task ethylic alcohol lên men thành acetic acid;
\task acetic acid ester hóa với ethylic alcohol.
\end{tasks}
\dapso{Xem phần lí thuyết.}
\end{tuluyen}

\begin{tuluyen}
Một chai rượu \SI{700}{\milli\liter}, độ cồn $35^\circ$. Tính thể tích ethylic alcohol nguyên chất trong chai.
\dapso{\SI{245}{\milli\liter}.}
\end{tuluyen}

\begin{tuluyen}
Có \SI{500}{\milli\liter} rượu $40^\circ$. Bỏ qua sự co thể tích khi pha, cần thêm bao nhiêu mililít nước để thu được rượu $25^\circ$?
\dapso{\SI{300}{\milli\liter} nước.}
\begin{loigiaiBT}
Thể tích ethylic alcohol ban đầu:
\[
V_{\ce{C2H5OH}}=500\times\frac{40}{100}=\SI{200}{\milli\liter}.
\]
Gọi thể tích dung dịch sau pha là $V$:
\[
\frac{200}{V}\times100=25\Rightarrow V=\SI{800}{\milli\liter}.
\]
\[
V_{\ce{H2O}}=800-500=\SI{300}{\milli\liter}.
\]

\textbf{Vậy,} cần thêm \SI{300}{\milli\liter} nước.
\end{loigiaiBT}
\end{tuluyen}

\begin{tuluyen}
Cho \SI{9.20}{\gram} ethylic alcohol tác dụng hoàn toàn với Na dư. Tính thể tích \ce{H2} thu được ở điều kiện có $V_m=\SI{24.8}{\liter\per\mole}$.
\dapso{\SI{2.48}{\liter}.}
\begin{loigiaiBT}
\[
n_{\ce{C2H5OH}}=\frac{9.20}{46}=\SI{0.200}{\mole}.
\]

\bangmol{7}
{\ce{2C2H5OH} & + & \ce{2Na} & \rightarrow & \ce{2C2H5ONa} & + & \ce{H2}}
{0.200 & & & & & & 0.100}

\[
V_{\ce{H2}}=0.100\times24.8=\SI{2.48}{\liter}.
\]

\textbf{Vậy,} thu được \SI{2.48}{\liter} \ce{H2}.
\end{loigiaiBT}
\end{tuluyen}

\begin{tuluyen}
Cho \SI{12.0}{\gram} acetic acid tác dụng hoàn toàn với đá vôi tinh khiết. Tính thể tích \ce{CO2} sinh ra ở điều kiện có $V_m=\SI{24.8}{\liter\per\mole}$.
\dapso{\SI{2.48}{\liter}.}
\begin{loigiaiBT}
\[
n_{\ce{CH3COOH}}=\frac{12.0}{60}=\SI{0.200}{\mole}.
\]

\bangmol{7}
{\ce{2CH3COOH} & + & \ce{CaCO3} & \rightarrow & \ce{(CH3COO)2Ca} & + & \ce{CO2 + H2O}}
{0.200 & & & & & & 0.100}

\[
V_{\ce{CO2}}=0.100\times24.8=\SI{2.48}{\liter}.
\]

\textbf{Vậy,} thu được \SI{2.48}{\liter} \ce{CO2}.
\end{loigiaiBT}
\end{tuluyen}

\begin{tuluyen}
Một chai giấm có thể tích \SI{750}{\milli\liter}, khối lượng riêng \SI{1.04}{\gram\per\milli\liter}, chứa $5.0\%$ acetic acid theo khối lượng. Tính khối lượng acetic acid trong chai.
\dapso{\SI{39.0}{\gram}.}
\end{tuluyen}

\begin{tuluyen}
Lên men glucose theo phản ứng \ce{C6H12O6 -> 2C2H5OH + 2CO2}.\par
Từ \SI{18.0}{\kilo\gram} glucose, hiệu suất lên men đạt $75\%$. Tính khối lượng ethylic alcohol thực tế thu được.
\dapso{\SI{6.90}{\kilo\gram}.}
\begin{loigiaiBT}
\[
n_{\ce{C6H12O6}}=\frac{18.0}{180}=\SI{0.100}{\kilo\mole}.
\]

\bangmol{5}
{\ce{C6H12O6} & \rightarrow & \ce{2C2H5OH} & + & \ce{2CO2}}
{0.100 & & 0.200 & &}

\[
m_{\ce{C2H5OH},\text{lt}}=0.200\times46=\SI{9.20}{\kilo\gram},
\]
\[
m_{\ce{C2H5OH}}=9.20\times0.75=\SI{6.90}{\kilo\gram}.
\]

\textbf{Vậy,} thu được \SI{6.90}{\kilo\gram} ethylic alcohol.
\end{loigiaiBT}
\end{tuluyen}

\begin{tuluyen}
Một mẫu giấm được pha loãng $20$ lần. Lấy \SI{10.00}{\milli\liter} dung dịch sau pha loãng cần \SI{4.00}{\milli\liter} dung dịch \ce{NaOH} \SI{0.100}{\mole\per\liter} để trung hòa. Tính nồng độ mol của acetic acid trong giấm ban đầu.
\dapso{\SI{0.800}{\mole\per\liter}.}
\end{tuluyen}

\subsection{Bài tập tổng hợp}

\begin{tuluyen}
Một chai giấm có thể tích \SI{1.00}{\liter}, khối lượng riêng \SI{1.045}{\gram\per\milli\liter}, chứa $4.0\%$ acetic acid theo khối lượng. Tính khối lượng acetic acid có trong chai.
\source{Phỏng theo Câu 5, đề tuyển sinh chuyên Hóa Lâm Đồng 2026.}
\dapso{\SI{41.8}{\gram}.}
\end{tuluyen}

\begin{tuluyen}
Trong thí nghiệm điều chế ethyl acetate, hỗn hợp ethylic alcohol và acetic acid được đun với một ít sulfuric acid đặc trong bình có ống sinh hàn. Sau phản ứng, hỗn hợp được làm nguội rồi thêm dung dịch \ce{NaCl} bão hòa trước khi tách lớp.
\begin{enumerate}
\item Viết phương trình phản ứng ester hóa.
\item Nêu tác dụng của ống sinh hàn.
\item Giải thích vì sao nước làm lạnh nên đi vào đầu dưới và ra ở đầu trên của ống sinh hàn.
\item Nêu vai trò của dung dịch \ce{NaCl} bão hòa trong bước tách sản phẩm.
\item Vì sao không thể coi phản ứng ester hóa là phản ứng xảy ra hoàn toàn?
\end{enumerate}
\source{Phỏng theo cấu trúc câu thí nghiệm ester hóa, đề tuyển sinh chuyên Hóa Lâm Đồng 2026.}
\dapso{Phản ứng thuận nghịch; sinh hàn hạn chế thất thoát chất dễ bay hơi; nước vào dưới để áo nước luôn đầy; NaCl bão hòa hỗ trợ tách ester khỏi pha nước.}
\end{tuluyen}

\begin{tuluyen}
Cho \SI{18.0}{\gram} acetic acid và \SI{18.4}{\gram} ethylic alcohol phản ứng ester hóa. Sau phản ứng thu được \SI{17.6}{\gram} ethyl acetate. Tính hiệu suất phản ứng.
\dapso{$66.7\%$.}
\begin{loigiaiBT}
\[
n_{\ce{CH3COOH}}=\frac{18.0}{60}=\SI{0.300}{\mole},\qquad
n_{\ce{C2H5OH}}=\frac{18.4}{46}=\SI{0.400}{\mole}.
\]

\bangmol{7}
{\ce{CH3COOH} & + & \ce{C2H5OH} & \rightleftharpoons & \ce{CH3COOC2H5} & + & \ce{H2O}}
{0.300 & & 0.400 & & 0.300 & &}

\[
m_{\text{ester,lt}}=0.300\times88=\SI{26.4}{\gram}.
\]
\[
H=\frac{17.6}{26.4}\times100\%\approx66.7\%.
\]

\textbf{Vậy,} hiệu suất phản ứng khoảng $66.7\%$.
\end{loigiaiBT}
\end{tuluyen}

\begin{tuluyen}
Một mẫu giấm được phân tích bằng chuẩn độ. Lấy \SI{10.00}{\milli\liter} giấm, pha loãng thành \SI{100.0}{\milli\liter}. Lấy \SI{10.00}{\milli\liter} dung dịch sau pha loãng, chuẩn độ bằng \ce{NaOH} \SI{0.100}{\mole\per\liter}. Ba lần chuẩn độ dùng $7.95$, $8.00$ và $8.05$ mL. Tính nồng độ mol acetic acid trong giấm ban đầu.
\source{Phỏng theo bài chuẩn độ giấm trong đề tuyển sinh chuyên Hóa Quảng Ninh 2026.}
\dapso{\SI{0.800}{\mole\per\liter}.}
\begin{loigiaiBT}
\[
\overline V=\frac{7.95+8.00+8.05}{3}=\SI{8.00}{\milli\liter}.
\]
\[
n_{\ce{NaOH}}=0.100\times8.00\times10^{-3}=\SI{8.00e-4}{\mole}.
\]

\bangmol{5}
{\ce{CH3COOH} & + & \ce{NaOH} & \rightarrow & \ce{CH3COONa + H2O}}
{8.00\times10^{-4} & & 8.00\times10^{-4} & &}

\[
n_{\ce{CH3COOH},\,100\,\mathrm{mL}}=8.00\times10^{-4}\times10=\SI{8.00e-3}{\mole}.
\]
\[
C_M=\frac{8.00\times10^{-3}}{10.00\times10^{-3}}=\SI{0.800}{\mole\per\liter}.
\]

\textbf{Vậy,} nồng độ mol của acetic acid trong giấm là \SI{0.800}{\mole\per\liter}.
\end{loigiaiBT}
\end{tuluyen}

\begin{tuluyen}
Trong phép chuẩn độ giấm bằng \ce{NaOH}, một học sinh quên tráng burette bằng dung dịch \ce{NaOH} trước khi nạp dung dịch chuẩn. Nêu chiều sai lệch của kết quả tính hàm lượng acetic acid và giải thích.
\source{Phỏng theo lệnh hỏi về sai số chuẩn độ trong đề tuyển sinh chuyên Hóa Quảng Ninh 2026.}
\dapso{Kết quả có xu hướng lớn hơn giá trị thật vì dung dịch \ce{NaOH} trong burette bị pha loãng.}
\end{tuluyen}
